\section{An exact operational separation of acquisition and receiver memory}
\label{sec:operationalmemory90}
The local profile law does not identify exact operational distances.
We now separate three strategy classes on one finite ensemble of ordered
measurement channels.  The separation holds in every input dimension,
throughout a depolarizing family, with exact optimal success probabilities.
The unknown interface still has only classical output.  The experiment
has a binary hypothesis decision, not a request to identify every member
of the ensemble.

\subsection{A finite family and its second moments}
Let $d\ge2$, let $S$ interchange the two factors of $\C^d\otimes\C^d$,
and put $\Pi_\pm=(I\pm S)/2$.  A weighted family of ordered orthonormal
bases consists of projections $P_y^b$, $1\le y\le d$, and positive
weights $w_b$ summing to one.  We require
The weights are real; rational or algebraic coordinates are not required
by the theorem.  We call this a weighted exact second-moment family,
not necessarily an unweighted projective design.
\begin{equation}\label{eq:basismoments90}
 \sum_b w_b P_y^b\otimes P_z^b=
 \begin{cases}
 (I+S)/(d(d+1)),&y=z,\\
 (dI-S)/(d(d^2-1)),&y\ne z.
 \end{cases}
\end{equation}
\begin{lemma}[Finite simultaneous second moments]\label{lem:finitedesign90}
In every dimension $d\ge2$ there is such a family with at most $d^6+1$
bases.  For $d=2$, the six ordered eigenbases of the three Pauli matrices,
with both orders and equal weights, satisfy \eqref{eq:basismoments90}.
\end{lemma}
\begin{proof}
Average the ordered coordinate basis over Haar measure on $U(d)$.
The second moment of a single column is supported on the symmetric
subspace, invariant under $U\otimes U$, and has trace one.  It is
therefore $2\Pi_+/[d(d+1)]$.  For distinct columns the invariant moment
has the form $aI+bS$; its trace is one and its contraction with $S$ is
zero, since the columns are orthogonal.  These two equations give the
second line of \eqref{eq:basismoments90}.  Equivalently, this line follows
by summing over the other $d-1$ columns and using the first moment $I/d$.
The vector of all $d^2$ Hermitian second moments lies in a real space of
dimension at most $d^6$.  Its Haar average belongs to the convex hull
of the compact set of these vectors.  Carath\'eodory's theorem gives a
convex combination of at most $d^6+1$ vectors; discard zero weights.
For the stated qubit family expand $P=(I\pm\sigma_a)/2$ and use
$S=(I\otimes I+\sum_{a=1}^3\sigma_a\otimes\sigma_a)/2$.
\end{proof}
This is a finite existence argument, not an efficient design construction.
Second-moment designs themselves are established objects \cite{GAE90}.

Fix $0\le t\le1$, and consider the ordered measurements
\begin{equation}\label{eq:gamefamily90}
 C_y=I/d,\qquad E_y^b(t)=(1-t)I/d+tP_y^b.
\end{equation}
Set
\begin{equation}\label{eq:gameprior90}
 L=2(d+1)-t^2,\qquad
 p_C=\frac{d+1-t^2}{L},\qquad p_E=\frac{d+1}{L}.
\end{equation}
With probability $p_C$ the device is $C$; otherwise a basis $b$ is chosen
with law $w_b$ and the device is $E^b(t)$.  This choice is made once.
Both calls use the same chosen, memoryless, classical-output device.
The task is to decide $C$ versus the alternative family.  In particular,
the second moments cannot be replaced by two independent first moments.
Every effect is strictly positive for $t<1$.
Taking the partial trace in \eqref{eq:basismoments90} gives
$\sum_b w_bE_y^b(t)=C_y$.  Thus the two hypotheses have identical
one-call channels, including on entangled inputs, and the one-call
optimum is $p_E$.  The two-call signal concerns a device chosen once;
it is absent if the alternative is independently redrawn at each call.

We compare: (i) complete first-call measurement followed by
measure-and-reprepare classical adaptation, as in
\cite[Definition~23]{OZNPQ89}; (ii) unit-reset protocols of
Definition~\ref{def:reset88}, with profile $(1,1)$ and arbitrary receiver
quantum memory and classical feedback; and (iii) all two-call adaptive
protocols.  Let their optimal Bayes success probabilities be
$P_{\rm ca}$, $P_{\rm reset}$ and $P_{\rm all}$, respectively.
Early stopping is allowed and can be implemented by a discarded fresh call.
The three closed optimization classes are precisely those of
Definition~\ref{def:classes91}.  Lemma~\ref{lem:resetnormal91} gives
the full branch normal form, and Proposition~\ref{prop:causalnormal91}
gives the physical causal normalization independently of this example.

\begin{theorem}[Exact two-call hierarchy]\label{thm:operationalhierarchy90}
The alternative basis index in this experiment is drawn once and
reused at both calls.  Independently redrawing it instead gives the
scalar product channel and no improvement over the larger prior.
For every finite family satisfying \eqref{eq:basismoments90}, every
$d\ge2$, and every $0\le t\le1$, the game \eqref{eq:gamefamily90}--\eqref{eq:gameprior90}
has the exact optimal values
\begin{equation}\label{eq:operationalhierarchy90}
\begin{split}
 P_{\rm ca}&=\frac{d+1}{L},\\
 P_{\rm reset}&=\frac{d+1}{L}+\frac{t^2(d-1)}{2dL},\\
 P_{\rm all}&=\frac{d+1+t^2}{L}.
\end{split}
\end{equation}
For $t>0$ both inequalities are strict.  At $t=0$ all three
values are $1/2$.  The priors sum to one and
$p_E-p_C=t^2/L$, so $p_E$ is the constant-decision baseline.
The middle value is attained
by independent maximally entangled input--reference pairs and a final
joint reference readout, without feedback.  Its upper bound allows all
unit-reset feedback and unrestricted receiver dimension.  The last
value is attained by a parallel antisymmetric two-input state, and its
upper bound allows every adaptive two-call tester.
\end{theorem}
The theorem separates a finite ensemble decision task, not an arbitrary
pair of fixed channels and not a complete memory-dimension hierarchy.
The middle class has hard resources $N=Q=2$; the last construction is
one acquisition group of size two, with square cost four.

\subsection{A swap-filter identity}
For a Hermitian matrix $X$, write $X_+$ for its positive part.
\begin{lemma}[Negative mass of a filtered swap]\label{lem:swapfilter90}
Let $A,B\succeq0$ be $d\times d$ matrices and
$K=(\sqrt A\otimes\sqrt B)S(\sqrt A\otimes\sqrt B)$.  Then
\begin{equation}\label{eq:swapfilter90}
 \tr(-K)_+=\frac12\left[
  \bigl(\tr\sqrt{\sqrt A B\sqrt A}\bigr)^2-\tr(AB)\right]
 \le\frac{d-1}{2d}\tr A\tr B.
\end{equation}
The same negative mass is obtained from $(U\otimes V)S(U\otimes V)^*$
when $U^*U=A$ and $V^*V=B$, even if the receiving spaces are larger.
If $\tr A\tr B>0$, equality in the upper bound holds if and only if
$A=(\tr A)I/d$ and $B=(\tr B)I/d$.
\end{lemma}
\begin{proof}
First suppose $A,B$ are positive definite.  The matrix $K$ is similar
to $(A\otimes B)S$.  Conjugation of the latter by $I\otimes B$ gives
$(AB\otimes I)S$.  The matrix $AB$ is similar to $\sqrt A B\sqrt A$;
let its positive eigenvalues be $\lambda_1,\ldots,\lambda_d$.
Conjugating by the tensor square of an eigenbasis matrix leaves $S$
unchanged.  The resulting matrix has eigenvalues $\lambda_i$ on the
one-dimensional spaces indexed by $(i,i)$, and
$\pm\sqrt{\lambda_i\lambda_j}$ on the two-dimensional spaces indexed
by $(i,j),(j,i)$, $i<j$.  Thus its negative mass is
$\sum_{i<j}\sqrt{\lambda_i\lambda_j}$, which gives the equality.
By Cauchy--Schwarz,
$\sum_i\lambda_i\ge d^{-1}(\sum_i\sqrt{\lambda_i})^2$.
Moreover, the polar decomposition and Hilbert--Schmidt Cauchy--Schwarz
give $\|\sqrt A\sqrt B\|_1\le\sqrt{\tr A\tr B}$.
These inequalities give the bound.  Approximation by positive definite
matrices proves the singular case.  The polar decompositions of $U,V$
identify their filtered swap with $K$ on its support, adding only zero
eigenvalues in the orthogonal complement.
For the equality assertion, put $f=\|\sqrt A\sqrt B\|_1$ and
$a=\tr A$, $b=\tr B$.  Equality in the stated bound, with $ab>0$,
forces equality in both $\tr(AB)\ge f^2/d$ and $f^2\le ab$.
The variational formula for the trace norm and equality in
Hilbert--Schmidt Cauchy--Schwarz imply $B=(b/a)A$: indeed a maximizing
unitary $U$ has $\sqrt B U=c\sqrt A$, and multiplication by its
adjoint gives this proportionality.  Equality in the first inequality
makes all eigenvalues of $\sqrt A B\sqrt A$ equal and positive.
Consequently $A=(a/d)I$ and $B=(b/d)I$.  These scalar matrices plainly
attain the bound.  This also excludes equality at nonzero singular
factors.
\end{proof}
The normalized equality statement and its singular case are expanded
in the proof of Lemma~\ref{lem:centeredswap91}: normalized fidelity one
forces equality of the two density matrices, and eigenvalue
Cauchy--Schwarz then forces all $d$ eigenvalues to be $1/d$.

\subsection{Payoff and complete normalization}
We use the fully transposed, unnormalized Choi representative
\begin{equation}\label{eq:choiconvention90}
 J^\sharp(E)=\left[(\operatorname{Id}\otimes\mathcal M_E)
          (|\Omega\rangle\langle\Omega|)\right]^{\mathsf T}
       =\sum_y E_y\otimes|y\rangle\langle y|,
 \qquad |\Omega\rangle=\sum_i|ii\rangle.
\end{equation}
Its output trace is $I$.  A simultaneous full transpose of the usual
tester representative makes its contraction with $J^\sharp$ the Born
probability.  Factor order is $I_1,Y_1,I_2,Y_2$; expressions grouping
$I_1,I_2$ use the canonical permutation of these factors.
Dephasing the output factors does not change any score.  Write $T_{0,yz}$
for the input block of the tester decision $C$.  The gain over always
guessing the alternative is
\begin{equation}\label{eq:payoff90}
 P-p_E=\sum_{y,z}\tr(T_{0,yz}W_{yz}),\qquad
 W_{yz}=p_C C_y\otimes C_z-p_E\sum_b w_b E_y^b(t)\otimes E_z^b(t).
\end{equation}
The averaged tensor in the second term is $(1-t^2)I/d^2$ plus $t^2$
times \eqref{eq:basismoments90}.  With $k=t^2/(dL)$, substitution gives
\begin{equation}\label{eq:payoffblocks90}
 W_{yy}=-kS,\qquad W_{yz}=-\frac{2k}{d-1}\Pi_-\quad(y\ne z).
\end{equation}
These operators are real in the coordinate basis; a full transpose
therefore does not alter the following congruences.

\begin{proof}[Proof of Theorem~\ref{thm:operationalhierarchy90}]
\emph{Classical adaptation.}
Each block of a measure-and-reprepare tester is a sum of positive
products $A\otimes B$.  Since
$\tr[(A\otimes B)S]=\tr(AB)\ge0$, equal labels contribute no positive
gain in \eqref{eq:payoff90}; different labels cannot contribute positively
either.  Hence $P_{\rm ca}\le p_E$, with equality by the constant decision.
This proof also covers limits of such testers.

\emph{All unit-reset protocols.}
Purify the first acquisition and write its state as
$\sum_i|i\rangle C_0|i\rangle$, with $\tr C_0^*C_0=1$.
After the first device label $y$, combine all intervening receiver
operations into an instrument with recorded outcome $h$.  For an upper
bound, retain its Stinespring environment in the receiver.  The first
conditional reference operator then has the form
$C_{yh}E_y^{\mathsf T}C_{yh}^*$, where
\[
 A_{yh}=C_{yh}^*C_{yh},\qquad
 \sum_h A_{yh}=C_0^*C_0\quad\hbox{for each }y.
\]
Retaining this environment only enlarges receiver processing and does
not import it into the second acquisition.  Conditional reset permits
a purification of that fresh acquisition with a map $D_{yh}$ satisfying
$B_{yh}=D_{yh}^*D_{yh}$ and $\tr B_{yh}=1$.
Its complete system is independent of the old receiver conditional on
$(y,h)$.  The two-reference payoff operator at record $(y,h,z)$ is
therefore
\[
 (C_{yh}\otimes D_{yh})W_{yz}^{\mathsf T}
                         (C_{yh}\otimes D_{yh})^*.
\]
For different labels it is nonpositive.  For equal labels, maximizing
any final effect bounded by the identity gives at most its positive
mass, which by Lemma~\ref{lem:swapfilter90} is
$k(d-1)\tr A_{yh}/(2d)$.  Summing over $y,h$ gives
\[
 P_{\rm reset}-p_E\le k(d-1)/2=\frac{t^2(d-1)}{2dL}.
\]
Public randomization can be included in the recorded history.  A normal
instrument with countably many outcomes is treated by trace-class
sums, or by finite truncations retaining a remainder outcome and then
taking the trace-norm limit.  The bound is independent of receiver
size.  Thus no finite-dimensional restriction on receiver storage has
entered the upper bound.
For $t>0$, the equality case of Lemma~\ref{lem:swapfilter90} also shows
that a nonrandomized protocol attaining this bound in the displayed
purified normal form must satisfy
$A_{yh}=(\tr A_{yh})I/d$ and $B_{yh}=I/d$ on every nonzero branch.
All deficits in the summed upper bound are nonnegative, so equality
forces their branchwise vanishing.  Summing over $h$ then gives
$C_0^*C_0=I/d$.  Maximal mixing is thus a consequence of equality,
not a restriction imposed on the competing protocols.  This statement
concerns the purified normal form, not uniqueness of its dilations
or of the final readout.

For attainment, prepare $|\Omega\rangle/\sqrt d$ independently on each
input and reference.  Declare $C$ precisely when the two device labels
agree and the references lie in $\Pi_-$.  The event probability under
$C$ is $a=(d-1)/(2d^2)$.  Under any $E^b(t)$ it is $(1-t^2)a$:
indeed $\tr E_y^b(t)=1$, $\tr E_y^b(t)^2=t^2+(1-t^2)/d$, and
$\tr[\Pi_-(E\otimes E)]=[(\tr E)^2-\tr E^2]/2$.
The factor $d^{-2}$ comes from the two normalized input--reference
pairs.  The resulting gain is
$a[p_C-p_E(1-t^2)]=t^2(d-1)/(2dL)$.

\emph{All adaptive protocols.}
The deterministic normalization of any two-call tester is
$T_0+T_1=\Xi\otimes I_{Y_2}$, where, after dephasing $Y_1$,
$\Xi$ has input blocks $\Xi_y\succeq0$ satisfying
\[
 \tr_{I_2}\Xi_y=\rho\quad(1\le y\le d),\qquad \tr\rho=1.
\]
Here is a direct realization of these constraints in the present
classical-output interface.  Purify the initial state as
$|\psi\rangle\in I_1\otimes R$, and let $\mathcal A_y$ be the
trace-preserving intervening map from $R$ to $I_2\otimes R'$ after
observing $y$.  Put
\[
 X_y=(\operatorname{Id}_{I_1}\otimes\mathcal A_y)
                 (|\psi\rangle\langle\psi|),\qquad
 \Xi_y=\tr_{R'}X_y,\qquad
 \rho=\tr_R|\psi\rangle\langle\psi|.
\]
Positivity and trace preservation give the displayed marginal equations.
If $F_{yz}$ is the final receiver effect for decision $C$, then
\[
 T_{0,yz}=\tr_{R'}\bigl[(I_{I_1I_2}\otimes\sqrt{F_{yz}})
           X_y(I_{I_1I_2}\otimes\sqrt{F_{yz}})\bigr].
\]
The Choi contraction \eqref{eq:choiconvention90} gives its Born
probability, and $0\preceq F_{yz}\preceq I$ gives
$0\preceq T_{0,yz}\preceq\Xi_y$.  This also proves positivity of the
complement, without an input-dimension factor from a normalized Choi
state.  In particular $\tr\Xi_y=1$.
Since $-kS\preceq kI$ and the other payoff blocks are nonpositive,
\[
 P_{\rm all}-p_E\le k\sum_y\tr\Xi_y=kd=t^2/L.
\]
Prepare any normalized antisymmetric pure state on the two inputs and
declare $C$ when their labels agree.  The event probability is $1/d$
under $C$, and $(1-t^2)/d$ under every $E^b(t)$, because
$(P_y^b\otimes P_y^b)\Pi_-=0$.  The gain is $t^2/L$.
This is a legitimate parallel two-call strategy and attains the
adaptive upper.  The case $t=0$ gives equal values $1/2$ directly.
\end{proof}

\begin{corollary}[Seven qubit devices]\label{cor:qubitgame90}
Take the scalar qubit measurement with prior $2/5$ and each of the six
ordered Pauli measurements with prior $1/10$.  With the device fixed
for both calls and the binary decision specified above, the exact
success probabilities are
\[
 P_{\rm ca}=3/5,\qquad P_{\rm reset}=13/20,\qquad P_{\rm all}=4/5.
\]
\end{corollary}
\begin{proof}
Use Lemma~\ref{lem:finitedesign90} and substitute $d=2,t=1$.
\end{proof}
The antisymmetric readout and comparison mechanism have antecedents in
measurement comparison \cite{ZHS90}; reductions between measurement and
state discrimination also have a substantial history \cite{SZ74}.
Here the assertion is the joint exact optimization over the three
specified classes, including unrestricted reset feedback, on one finite
classical-output ensemble and its full-rank deformation.

\subsection{Quantitative stability of the operational gap}
\begin{corollary}[A stable score separation]\label{cor:robustgame90}
Suppose every device in the finite game is replaced by a classical-output
measurement channel at diamond distance at most $\delta$ from its
specified channel.  Keep the priors and the two-call resource classes
fixed.  Each of their optimal success probabilities changes by at most
$\delta$.  If the displayed reset implementation also has total hard
pathwise preparation defect at most $\varepsilon$ in the sense of
Proposition~\ref{prop:approxreset89}, its score exceeds the optimal
perturbed classical-adaptation score whenever
\begin{equation}\label{eq:robustgap90}
 2\delta+\varepsilon/2<\frac{t^2(d-1)}{2dL}.
\end{equation}
\end{corollary}
\begin{proof}
A two-call hybrid changes each normalized output in unhalved trace norm
by at most $2\delta$.  A fixed decision event changes by at most
$\delta$, and averaging over the unchanged priors preserves this bound.
Taking suprema over a fixed strategy class gives the same two-sided
bound for its optimum.  The preparation certificate changes each
implemented output by at most $\varepsilon$, hence its decision
probability by at most $\varepsilon/2$.  Subtract the classical upper
from the reset lower in \eqref{eq:operationalhierarchy90}.
\end{proof}
These are conditional norm estimates, not physical calibration claims.
In the preceding corollary the same $\delta$ bounds the scalar channel
and each entire basis measurement channel individually.  Priors and the
once-drawn-index law stay fixed; index-dependent channel perturbations
are permitted.  Corollary~\ref{cor:gamestability91} additionally treats
prior and law perturbations, with the unhalved norm factors explicit.
They make the relative scale needed to preserve a small score gap explicit.

\subsection{Hard budget conventions beyond the normalized domain}
\begin{proposition}[Normalization of arbitrary hard thresholds]\label{prop:budgetdomain90}
For positive integer reservations, $N_\pi\le Q_\pi\le N_\pi^2$.
Given real hard bounds $n,q\ge1$, put
\[
 M=\min\{\lfloor n\rfloor,\lfloor q\rfloor\},\qquad
 R=\min\{\lfloor q\rfloor,M^2\}.
\]
The class $\{\pi:N_\pi\le n,\ Q_\pi\le q\}$ equals
$\{\pi:N_\pi\le M,\ Q_\pi\le R\}$, and $M\le R\le M^2$.
Consequently Theorem~\ref{thm:quadraticbudget89} applies with $(M,R)$.
\end{proposition}
\begin{proof}
On each path $\sum n(h)\le\sum n(h)^2\le(\sum n(h))^2$.
Taking hard suprema proves the first assertion.  Bounded integer costs
have integer suprema, so both constraints can be floored.  The first
inequality reduces the call bound to $M$ and the last reduces the
quadratic bound to $R$.  Conversely the reduced bounds imply the original
ones.  Finally $M\le\lfloor q\rfloor$ and $M\le M^2$.
\end{proof}
The normalization $N\le Q$ allows the $N$-singleton lower construction;
it does not say that every admissible individual policy has $Q_\pi\ge N$.
If a hard threshold is below one, only the zero-call class is available.
Fractional-duration reservations and expected-cost constraints are
different models.  None is substituted for a hard integer reservation here.
